\section{Participant Demographics}
\label{sec:eval-demographics}

This section presents the participant demographics.
The study, set out in full in Chapter~\ref{ch:methodology}, recruited $N = 73$ participants who self-reported
as ROPA novices for a within-subjects experiment in which every
participant used all three ROPAgen interaction modes (Form, Ask, and
Chat; see Section~\ref{sec:ropagen}).
Of the $73$ recruited participants, four screened records were
excluded for incomplete or invalid participation, yielding the matched
within-subjects analysis sample of $n = 69$ on which every
user-study result in this chapter rests.
Mode order was counterbalanced
across a six-group Balanced Latin Square so that learning, order, and
fatigue effects could be distributed evenly across interaction modes rather than
absorbed into any one of them. The task was designed around a single standardized GDPR scenario, the
APOR~GmbH employee and parking-management case at a fictional Ulm
software company, as described in Section~\ref{ch:methodology}\footnote{Full scenario described in \ref{app:scenario}}.
After interacting with each mode, participants completed the \textit{SUS}, the \textit{NASA R-TLX}, and the \textit{Perceived AI
Support} items, which are evaluated in the remaining sections of this chapter;
at the end of the session, they ranked the three modes
by overall preference and supplied the demographic information
profiled below. The demographic variables are reported descriptively only on the
matched within-subjects sample of $n = 69$ participants.\footnote{Descriptive
statistics for session duration and prior expertise, along with
counts and inferential tests for the categorical demographic
variables, are reported in Appendix~\ref{app:demographics}.}

The mean session duration was $1866$~s ($\approx 31$~min) per
participant across the full session, see Figure~\ref{fig:demographics-duration}.
Figure~\ref{fig:demographics-expertise} shows the reported ROPA familiarity with a mean of $1.45$.
Chatbot familiarity is correspondingly high, the mean being $4.36$,
while GDPR knowledge sits near the scale midpoint (mean $2.81$).
Figure~\ref{fig:demographics-expertise} makes the contrast vivid:
chatbot familiarity concentrates at the high end, ROPA familiarity
at the low end, and GDPR knowledge spreads broadly across the middle.
In addition, the sample is dominated by Bachelor-level students
(Figure~\ref{fig:demographics-degree}) studying or working in IT or
Computer Science (Figure~\ref{fig:demographics-field}).
Of the participants who reported gender (Figure~\ref{fig:demographics-gender}), $62.1\,\%$
identified as male.
The underlying age distribution (Figure~\ref{fig:demographics-age}) clusters tightly around
$24$--$26$ years, with a thin right tail of older participants.

This profile (IT-leaning, chatbot-fluent, ROPA-novice)
frames the results reported in the remainder of this chapter and
in Chapter~\ref{ch:docquality}, and is taken up again in the
Discussion chapter when generalizability is addressed.

\begin{figure}[!htbp]
\centering
\begin{subfigure}[t]{0.45\linewidth}
  \centering
  \includegraphics[width=\linewidth]{demographics_duration.png}
  \caption{Session duration (seconds).}
  \label{fig:demographics-duration}
\end{subfigure}\hfill
\begin{subfigure}[t]{0.45\linewidth}
  \centering
  \includegraphics[width=\linewidth]{demographics_expertise.png}
  \caption{Self-reported prior expertise (GDPR knowledge, ROPA
  familiarity, chatbot familiarity) on a 5-point Likert scale
  ($1$ = no familiarity, $5$ = expert).}
  \label{fig:demographics-expertise}
\end{subfigure}\\[1ex]
\begin{subfigure}[t]{0.45\linewidth}
  \centering
  \includegraphics[width=\linewidth]{demographics_age.png}
  \caption{Age distribution.}
  \label{fig:demographics-age}
\end{subfigure}\hfill
\begin{subfigure}[t]{0.45\linewidth}
  \centering
  \includegraphics[width=\linewidth]{demographics_gender.png}
  \caption{Self-reported gender ($n = 66$; not all participants
  answered this item).}
  \label{fig:demographics-gender}
\end{subfigure}\\[1ex]
\begin{subfigure}[t]{0.45\linewidth}
  \centering
  \includegraphics[width=\linewidth]{demographics_degree.png}
  \caption{Highest completed degree ($n = 68$; not all
  participants answered this item).}
  \label{fig:demographics-degree}
\end{subfigure}\hfill
\begin{subfigure}[t]{0.45\linewidth}
  \centering
  \includegraphics[width=\linewidth]{demographics_profession.png}
  \caption{Current professional status.}
  \label{fig:demographics-profession}
\end{subfigure}\\[1ex]
\begin{subfigure}[t]{0.45\linewidth}
  \centering
  \includegraphics[width=\linewidth]{demographics_field.png}
  \caption{Field of study or work.}
  \label{fig:demographics-field}
\end{subfigure}
\caption{Participant demographics of the $n = 69$ analysis sample:
session duration, self-reported prior expertise on a 5-point
Likert scale, and five categorical variables (age, gender,
highest completed degree, professional status, and field of study
or work).}
\label{fig:demographics-all}
\end{figure}

\FloatBarrier
